\documentclass[12pt,a4paper]{article}
\usepackage[utf8]{inputenc}
\usepackage[russian]{babel}
\usepackage[T1,T2A]{fontenc}
\usepackage[margin=2.5cm]{geometry}
\usepackage{hyperref}
\usepackage{amsmath}
\usepackage{amssymb}

\setlength{\parindent}{0pt}
\setlength{\parskip}{0.7em}

\begin{document}

\begin{center}
    \LARGE\textbf{Текст выступления перед комиссией}\\[1em]
    \Large Прайсинг деривативов с помощью нейронных сетей\\[1em]
    \large \textbf{Бирюков Никита Алексеевич, МФТИ}\\
    \textit{Научный руководитель: Андреев Михаил Владимирович}
\end{center}

\vspace{1.5em}

\section*{Слайд 1 — Титульный лист}

Уважаемые члены комиссии, уважаемый научный руководитель!

Меня зовут Бирюков Никита Алексеевич, и я представляю вашему вниманию дипломную работу на тему «Прайсинг деривативов с помощью нейронных сетей».

Моя работа посвящена созданию универсального фреймворка для нейросетевого прайсинга и калибровки различных моделей. Реализация суррогата стохастической модели волатильности Хестона выступает как первый шаг (Proof of Concept), чтобы оценить применимость такого пайплайна в реальном продакшене.

\vspace{1em}\hrule\vspace{1em}

\section*{Слайд 2 — Постановка задачи}

Позвольте начать с мотивации. Классическая калибровка модели Хестона — это итеративная процедура: на каждом шаге оптимизатора необходимо вычислять характеристическую функцию, оценивать 88 цен опционов через преобразование Фурье по методу Карра–Мадана и считать градиенты численно. Всё это повторяется до 10 000 раз — и занимает от 10 до 120 секунд на одну калибровку.

Для реал-тайм трейдинга это неприемлемо: жёсткий дедлайн — менее 150 миллисекунд. Кроме того, численные градиенты не позволяют точно вычислять Греки второго порядка.

\textbf{Главная цель работы} — разработать универсальный ML-пайплайн для решения этих проблем на масштабе. В качестве первого этапа и доказательства концепции мы рассматриваем модель Хестона, реализуя для неё трёхфазный пайплайн: обучение MLP-суррогата, байесовскую оценку неопределённости и LSTM-прогноз динамики параметров.

\vspace{1em}\hrule\vspace{1em}

\section*{Слайд 3 — Обзор ML-подходов}

В литературе существует несколько направлений применения нейронных сетей к задачам прайсинга деривативов.

Метод Deep BSDE (Han et al., 2018) решает нелинейные уравнения в частных производных, но не строит полную поверхность волатильности и не предоставляет аналитического якобиана. DGM от Sirignano требует ресурсов уровня V100/A100 и также не решает задачу быстрой калибровки. Deep Hedging решает принципиально другую задачу — оптимизацию хеджирования с транзакционными издержками.

Нас интересует направление Deep Calibration (Horvath et al., 2019): сеть обучается один раз в оффлайн-режиме, после чего инференс занимает миллисекунды. Автоматическое дифференцирование через autograd даёт точный якобиан, необходимый для L-BFGS-B. При этом требования к железу минимальны — RTX 3060 или CPU.

\vspace{1em}\hrule\vspace{1em}

\section*{Слайд 4 — Архитектура системы}

Система состоит из трёх последовательных фаз.

\textbf{Фаза 1} — MLP-суррогат: по вектору параметров Хестона ($\kappa, \theta, \sigma, \rho, v_0$) сеть предсказывает всю поверхность волатильности из 88 точек. На входе — MinMax-нормализация, на выходе — обратная стандартизация.

\textbf{Фаза 2} — MC Dropout: тот же MLP, но с оставленным активным Dropout на инференсе. Сто стохастических проходов дают доверительный интервал $\pm 2\sigma$ для каждой точки поверхности.

\textbf{Фаза 3} — LSTM-динамика: на основе истории поверхностей за 10 дней сеть прогнозирует параметры Хестона на следующий день, обеспечивая автоматическое соблюдение условия Феллера.

Фазы связаны: поверхности из Фазы 1 являются обучающими данными для Фазы 3.

\vspace{1em}\hrule\vspace{1em}

\section*{Слайд 5 — MLP-суррогат}

Перейдём к архитектуре суррогатной модели. Сеть имеет 4 скрытых слоя по 30 нейронов, итого 5 698 обучаемых параметров — это компактная модель, которая обучается за около 2 минут на RTX 3060.

Ключевое решение — выбор пространства предсказания. Сеть предсказывает не подразумеваемую волатильность напрямую, а суммарную дисперсию $W(K,T) = \sigma_{\text{IV}}^2(K,T) \cdot T$. Это пространство линеаризует уравнение Дюпира и делает естественным критерий Карра–Мадана для проверки отсутствия арбитража.

Вместо ReLU используется активация ELU. Причина принципиальная: норма Гессиана для ELU равна 1.039, тогда как для ReLU она строго равна нулю. Это означает, что Греки второго порядка, вычисляемые через autograd, для ReLU-сети будут тождественно нулевыми — что делает их бессмысленными для задач хеджирования.

\vspace{1em}\hrule\vspace{1em}

\section*{Слайд 6 — Результаты калибровки}

На слайде представлены результаты калибровки на тестовом примере. Синяя поверхность — целевая рыночная подразумеваемая волатильность, оранжевая — откалиброванная нейросетью.

Визуально поверхности практически совпадают. В численном выражении: RMSE в пространстве $W$ составляет \textbf{0.0876}, время одной калибровки — \textbf{47–135 миллисекунд}, нарушений условия арбитража — \textbf{ноль из 77 тестовых примеров}, условие Феллера выполнено во всех случаях.

Параметры восстанавливаются с высокой точностью: наибольшая ошибка по $\rho$ — 0.073, что обусловлено известной нечувствительностью поверхности к этому параметру при малых $v_0$.

\vspace{1em}\hrule\vspace{1em}

\section*{Слайд 7 — Производительность}

Бенчмарк скорости показывает кардинальную разницу. Классический подход FFT плюс Левенберг–Марквардт занимает на CPU 10–120 секунд — на графике это полоса длиной 12 единиц. Нейросетевой суррогат с L-BFGS-B: \textbf{0.091 секунды}, то есть \textbf{ускорение в 160 раз}.

На GPU разрыв ещё больше: классика — около 5 секунд, суррогат — менее 10 миллисекунд.

Причины ускорения: вместо 10 000 итераций FFT выполняется один forward-pass MLP; точный якобиан через autograd позволяет L-BFGS-B сходиться за единицы итераций, а не сотни.

Это открывает применение в high-frequency трейдинге, портфельной переоценке в реальном времени и батч-калибровке тысяч инструментов.

\vspace{1em}\hrule\vspace{1em}

\section*{Слайд 8 — Байесовская неопределённость (MC Dropout)}

Точечная оценка параметров не говорит, насколько уверена модель в своём предсказании. Фаза 2 решает эту задачу методом MC Dropout в интерпретации Гала и Ганрамани (2016).

Идея: оставить Dropout активным на инференсе и выполнить $N = 100$ стохастических forward-pass. Среднее даёт оценку поверхности, стандартное отклонение — меру неопределённости. Это является приближённым байесовским выводом над весами сети.

Результаты для $N = 100$: максимальное $\sigma_{\text{MC}}$ в ATM-точке при $T = 1$ составляет менее $5 \times 10^{-4}$, ширина полос $\pm 2\sigma$ — менее \textbf{0.1\% подразумеваемой волатильности}. Накладные расходы — примерно $100\times$ по сравнению с одним forward-pass, что всё равно укладывается в требование 150 мс.

Узкие полосы подтверждают отсутствие переобучения и хорошую генерализацию модели.

\vspace{1em}\hrule\vspace{1em}

\section*{Слайд 9 — LSTM-динамика параметров}

Третья фаза — прогноз динамики параметров Хестона во времени. Задача: по истории поверхностей $W_t$ за 10 дней предсказать вектор параметров $\hat{\boldsymbol\theta}_{t+1}$ на следующий день.

Архитектура: двухслойный LSTM с 64 скрытыми нейронами, LayerNorm после LSTM, 73 157 параметров. Обучение на AdamW с CosineAnnealingLR и ранней остановкой с patience 20.

Обучающий датасет генерировался следующим образом: 600 траекторий параметров по процессам Орнштейна–Уленбека, откалиброванным к динамике SPX; на каждом шаге — вызов MLP-суррогата для получения поверхности; итого \textbf{102 000 скользящих окон} $(W_{t-9}, \ldots, W_t) \to \boldsymbol\theta_{t+1}$.

Разбивка 80/10/10 делалась по границам траекторий, что исключает утечку данных между обучением и тестом.

Метки нормализованы по Z-score, на инференсе применяется обратное преобразование с жёсткой привязкой к физически допустимой области — условие Феллера гарантировано на каждом вызове.

\vspace{1em}\hrule\vspace{1em}

\section*{Слайд 10 — Результаты LSTM}

Нормированный RMSE на тестовой выборке из 10 200 последовательностей составил \textbf{0.336}, нормированный MSE — \textbf{0.113}.

Наиболее точно прогнозируются параметры с медленной динамикой: $\theta$ (RMSE = 0.0035) и $v_0$ (0.0052) — они следуют OU-процессу с большой инерцией. Параметр $\kappa$ — скорость возврата к среднему — прогнозируется хуже всего (RMSE = 0.127), что объясняется его высокой волатильностью и слабой идентификацией из поверхности волатильности.

Это согласуется с известными результатами в литературе о том, что $\kappa$ является наименее устойчивым параметром при калибровке модели Хестона.

\vspace{1em}\hrule\vspace{1em}

\section*{Слайд 11 — Демонстрация}

Все три фазы реализованы в интерактивном Streamlit-приложении, доступном по адресу \textbf{hestondl.streamlit.app}.

Вкладка «Демо калибровки» позволяет задать параметры ползунками, сгенерировать поверхность одним кликом, запустить калибровку с 3D-сравнением результата, оценить неопределённость методом MC Dropout и получить LSTM-прогноз на следующий день.

Тестовое покрытие проекта составляет 29 автоматических pytest-тестов: 7 тестов для Фазы 2 и 22 — для Фазы 3. Все тесты проходят на 100\%.

\vspace{1em}\hrule\vspace{1em}

\section*{Слайд 12 — Выводы и дальнейшие направления}

Подведём итоги. В работе был разработан масштабируемый трёхфазный пайплайн для нейросетевого прайсинга. На примере модели Хестона мы доказали, что этот фреймворк готов к продакшену:

\textbf{Фаза 1} — MLP-суррогат: RMSE = 0.0876, время калибровки 47–135 мс, ускорение в 160 раз по сравнению с классикой, ноль нарушений арбитражных ограничений.

\textbf{Фаза 2} — MC Dropout: байесовские доверительные полосы шириной менее 0.1\% IV без изменения архитектуры сети.

\textbf{Фаза 3} — LSTM: нормированный RMSE 0.336 при 102 000 обучающих окнах, гарантированное выполнение условия Феллера.

Поскольку пайплайн универсален, среди главных направлений развития — расширение фреймворка на другие модели, в том числе с грубой волатильностью (rBergomi, rHeston), калибровка на реальных исторических данных SPX, использование среды для deep hedging и полноценное развёртывание через REST API в облаке.

\vspace{1.5em}
\begin{center}
    \textit{Благодарю за внимание и готов ответить на ваши вопросы.}
\end{center}

\end{document}
